\documentclass[aps,prd,twocolumn,superscriptaddress,nofootinbib]{revtex4-2}

\usepackage{amsmath}
\usepackage{amssymb}
\usepackage{amsthm}
\usepackage{booktabs}
\usepackage{tabularx}
\usepackage{xcolor}
\usepackage{mathrsfs}

\newtheorem{theorem}{Theorem}[section]
\newtheorem{proposition}[theorem]{Proposition}
\newtheorem{lemma}[theorem]{Lemma}
\newtheorem{corollary}[theorem]{Corollary}
\newtheorem{conjecture}[theorem]{Conjecture}
\newtheorem{axiom}[theorem]{Axiom}
\theoremstyle{definition}
\newtheorem{definition}[theorem]{Definition}
\theoremstyle{remark}
\newtheorem{remark}[theorem]{Remark}
\newtheorem{claim}[theorem]{Claim}

\DeclareMathOperator{\poly}{poly}
\DeclareMathOperator{\SIZE}{SIZE}
\DeclareMathOperator{\SAT}{SAT}
\newcommand{\NP}{\mathrm{NP}}
\newcommand{\PP}{\mathrm{P}}
\newcommand{\BPP}{\mathrm{BPP}}
\newcommand{\coNP}{\mathrm{coNP}}

\usepackage{hyperref}
\hypersetup{
    pdftitle={P vs NP as an Entropic No-Go Theorem: Algorithmic Positivity, Kolmogorov Obstruction, and the Impossibility of Witness Compression},
    pdfauthor={Lukas Geiger},
    colorlinks=true,
    linkcolor=black,
    urlcolor=blue,
    citecolor=black
}

\begin{document}

% ===================================================================
% TITLE PAGE
% ===================================================================

\title{P vs NP as an Entropic No-Go Theorem:\\Algorithmic Positivity, Kolmogorov Obstruction,\\and the Impossibility of Witness Compression}

\author{Lukas Geiger}
\email{https://orcid.org/0009-0005-7296-1534}
\affiliation{Independent Researcher, Bernau im Schwarzwald, Germany}

\date{March 2026}

\begin{abstract}
We propose a reformulation of the $\PP$ vs $\NP$ problem as an \textit{entropic No-Go theorem}: $\PP \neq \NP$ is equivalent to the structural impossibility of compressing NP-witnesses to algorithmic entropy below their intrinsic Kolmogorov complexity. This paper does not prove $\PP \neq \NP$; rather, it provides an information-theoretic reformulation that identifies the precise structure of what must be shown and why classical barrier techniques are not obstructions to this approach. We introduce the concept of \textit{algorithmic positivity} -- an NP language is algorithmically positive if it admits an entropically convex, polynomially scaling witness-producing function -- and show that the \textit{Entropic Separation Conjecture} (no NP-complete language is algorithmically positive) is \textit{equivalent} to $\PP \neq \NP$. We conduct a systematic barrier analysis, arguing that Kolmogorov-complexity-based entropy arguments are not naturally subject to the three classical barriers (relativization, natural proofs, algebrization). The key structural insight is the \textit{witness entropy gap}: $\PP = \NP$ would imply $K^t(w \mid x) = O(\log|x|)$ for all NP-witnesses, while $\PP \neq \NP$ implies the existence of instance families with $K^t(w \mid x) \geq |w| - O(\log n)$ for every polynomial time bound $t$. We identify the precise remaining challenge: converting this characterization into a proof requires a \textit{uniformity bridge} -- a new connection between algorithmic information theory and uniform complexity theory.
\end{abstract}

\keywords{P vs NP, Kolmogorov complexity, algorithmic entropy, No-Go theorem, barrier analysis, witness compression, computational complexity}

\thanks{AI tools (Claude by Anthropic) were used for mathematical formalization and text generation. All content and responsibility for correctness lie solely with the author.}

\maketitle


% ===================================================================
% 1. INTRODUCTION
% ===================================================================
\section{Introduction}
\label{sec:intro}

\subsection{The problem}

The $\PP$ vs $\NP$ question asks whether every problem whose solutions can be \textit{verified} in polynomial time can also be \textit{solved} in polynomial time. Formally:

\begin{definition}
$\PP$ is the class of languages decidable in polynomial time. $\NP$ is the class of languages for which membership has a polynomial-time verifiable witness~\cite{Arora2009}. $\PP = \NP$ iff every language in $\NP$ is also in $\PP$. By the Cook-Levin theorem~\cite{Cook1971}, this is equivalent to the satisfiability problem (SAT) being solvable in polynomial time.
\end{definition}

Despite decades of effort, the problem remains open. This paper argues that the difficulty is not technical but \textit{conceptual}: previous approaches attempt to prove lower bounds via \textit{constructive} methods (exhibiting specific hard instances, explicit circuit separations, or algebraic obstructions), which are systematically blocked by three barriers. We propose an alternative: prove $\PP \neq \NP$ via a \textit{structural entropy argument} that identifies the impossibility of witness compression as the fundamental obstruction.

\subsection{The three barriers}

Three meta-theorems limit the scope of known proof techniques:

\textbf{1. Relativization} \cite{BakerGillSolovay1975}. There exist oracles $A, B$ with $\PP^A = \NP^A$ and $\PP^B \neq \NP^B$. Therefore, any proof of $\PP \neq \NP$ must use techniques that do not relativize -- it must exploit the internal structure of computation, not just input-output behavior.

\textbf{2. Natural proofs} \cite{RazborovRudich1997}. If strong pseudorandom generators exist (a widely believed assumption), then no ``natural'' proof method -- one that is both \textit{constructive} (recognizes hard functions efficiently) and \textit{large} (applies to a dense set of functions) -- can prove superpolynomial circuit lower bounds against $\PP/\poly$.

\textbf{3. Algebrization} \cite{AaronsonWigderson2009}. Even non-relativizing techniques based on arithmetization (as in $\mathrm{IP} = \mathrm{PSPACE}$) cannot separate $\PP$ from $\NP$, because there exist algebraic extensions of oracles in which both $\PP = \NP$ and $\PP \neq \NP$ hold.

\subsection{The positivity alternative}

We observe that major breakthroughs across mathematics share a common structure: they succeed not by constructing the desired object, but by showing its non-existence would violate a \textit{positivity condition}:

\begin{center}
\begin{tabular}{lll}
\hline
Problem & Positive structure & Status \\
\hline
Weil conjectures & Frobenius eigenvalues & Proven \\
BSD (rank $\leq 1$) & N\'eron-Tate height & Proven \\
Hodge conjecture & Hodge-Riemann form & Open \\
RH & Spectral zeta flow & \cite{GeigerRH2026} \\
\textbf{P vs NP} & \textbf{Kolmogorov complexity} & \textbf{This paper} \\
Broader program & Free energy framework & \cite{GeigerFST2026} \\
\hline
\end{tabular}
\end{center}

The key insight: \textit{Kolmogorov complexity is the natural positivity structure for computational complexity}. It is non-negative ($K(x) \geq 0$), machine-invariant (up to additive constants), and fundamentally non-constructive (not computable) -- precisely the properties needed to bypass the three barriers.


% ===================================================================
% 2. BARRIER ANALYSIS
% ===================================================================
\section{Systematic Barrier Analysis}
\label{sec:barriers}

\subsection{Why constructive methods fail}

All three barriers share a common cause: they block proofs that treat computational problems as \textit{black boxes} or that rely on \textit{statistical} properties of Boolean functions.

\begin{itemize}
\item \textbf{Relativization} blocks proofs that treat the computation as input-output behavior (oracle access), ignoring internal structure.
\item \textbf{Natural proofs} block proofs that identify hard functions via efficiently testable, statistically dense properties.
\item \textbf{Algebrization} blocks proofs that use only the polynomial structure of arithmetized computations.
\end{itemize}

The common thread: all three barriers apply to \textit{model-dependent} arguments -- arguments that work within a specific computational model (oracle machines, circuits, algebraic extensions) rather than targeting the \textit{intrinsic information content} of the computation.

\subsection{The specification for a valid proof}

From the barrier analysis, we derive necessary conditions for any proof of $\PP \neq \NP$:

\begin{enumerate}
\item[\textbf{(S1)}] \textbf{Non-relativizing:} The argument must exploit internal computational structure, not just oracle behavior.
\item[\textbf{(S2)}] \textbf{Non-natural:} The argument must not define a ``largeness'' property that is efficiently computable. It must target individual objects or non-constructive properties.
\item[\textbf{(S3)}] \textbf{Non-algebrizing:} The argument must go beyond polynomial identity testing and arithmetization.
\item[\textbf{(S4)}] \textbf{Model-independent:} The argument must not depend on a specific computational model (circuits, Turing machines, etc.) but on an invariant structural property.
\item[\textbf{(S5)}] \textbf{Robust:} The argument must not break when moving from restricted models (monotone circuits, $\mathrm{AC}^0$) to general computation.
\end{enumerate}

We now show that algorithmic entropy satisfies all five specifications.


% ===================================================================
% 3. ALGORITHMIC ENTROPY
% ===================================================================
\section{Algorithmic Entropy as Positivity Structure}
\label{sec:entropy}

\subsection{Kolmogorov complexity}

Let $U$ be a fixed universal prefix-free Turing machine. The \textit{Kolmogorov complexity} (algorithmic entropy) of a finite binary string $x$ is:
\begin{equation}
K(x) := \min\{|p| : U(p) = x\}\,.
\label{eq:kolmogorov}
\end{equation}

The \textit{conditional} Kolmogorov complexity of $x$ given $y$ is:
\begin{equation}
K(x \mid y) := \min\{|p| : U(p, y) = x\}\,.
\label{eq:conditional_K}
\end{equation}

\begin{theorem}[Invariance {\cite{Kolmogorov1965,Solomonoff1964,Chaitin1966}}]
\label{thm:invariance}
For any two universal prefix-free machines $U, U'$:
\begin{equation}
|K_U(x) - K_{U'}(x)| \leq c_{U,U'}\,,
\end{equation}
where $c_{U,U'}$ is a constant independent of $x$.
\end{theorem}

\begin{theorem}[Incomputability]
\label{thm:incomputable}
$K(x)$ is not computable. More precisely: for every total computable function $f$, there exists $x$ with $K(x) > f(x)$.
\end{theorem}

\begin{theorem}[Incompressibility]
\label{thm:incompressible}
For every $n$, at least $2^n - 2^{n-c+1} + 1$ strings of length $n$ satisfy $K(x) \geq n - c$.
\end{theorem}

\begin{theorem}[Symmetry of Information {\cite{LiVitanyi2008}}]
\label{thm:symmetry}
For all strings $x, y$:
\begin{equation}
K(x, y) = K(x) + K(y \mid x, K(x)) + O(\log(K(x, y)))\,.
\end{equation}
This implies that conditional complexity is well-behaved: the information in the pair $(x, y)$ decomposes additively (up to logarithmic terms) into the information in $x$ plus the additional information in $y$ given $x$.
\end{theorem}

\subsection{Why Kolmogorov complexity bypasses the barriers}

\begin{claim}[Barrier immunity]
\label{claim:barrier_immune}
Arguments based on Kolmogorov complexity plausibly satisfy specifications (S1)--(S5). We present heuristic arguments for each; a fully rigorous treatment is an open research question (see Limitations, Section~\ref{sec:discussion}).
\end{claim}

\begin{enumerate}
\item[\textbf{(S1)}] \textbf{Non-relativizing.} $K(x)$ is defined without reference to any oracle -- it is an absolute, fixed quantity for each string $x$. The witness entropy gap (Theorem~\ref{thm:high_entropy}) uses $K^t(w \mid x)$ (resource-bounded conditional complexity). We note an important subtlety: the invariance theorem for resource-bounded $K^t$ is weaker than for unbounded $K$ -- changing the universal machine can introduce a \textit{polynomial} slowdown factor rather than just an additive constant~\cite{LiVitanyi2008}. However, the argument remains non-relativizing because $K^t$ is defined over a fixed computational model without oracle access. A relativizing proof would need to hold for all oracles $A$, but our argument targets oracle-independent quantities: the gap between $O(\log n)$ and $\Omega(|w|)$ is superpolynomial and thus robust under polynomial slowdown factors.

\item[\textbf{(S2)}] \textbf{Non-natural.} $K(x)$ is not computable (Theorem~\ref{thm:incomputable}), so no efficient procedure can test whether a specific function has ``high Kolmogorov complexity.'' This directly violates the ``constructiveness'' requirement of natural proofs.

\item[\textbf{(S3)}] \textbf{Non-algebrizing.} Algebrization applies to proof techniques that use only the polynomial structure of arithmetized computations. Arguments based on $K^t$ operate over the full power of Turing machine computation, not polynomial identity testing. While it remains a research question whether $K^t$-based arguments are formally non-algebrizing in the technical sense of~\cite{AaronsonWigderson2009}, the entropy framework is not naturally expressible as an algebraic identity, suggesting it escapes this barrier.

\item[\textbf{(S4)}] \textbf{Model-independent.} The invariance theorem ensures that $K(x)$ is the same (up to $O(1)$) for all universal machines -- Turing machines, circuits, $\lambda$-calculus, etc.

\item[\textbf{(S5)}] \textbf{Robust.} $K(x)$ does not depend on the computational model used to define it.
\end{enumerate}


% ===================================================================
% 4. THE WITNESS ENTROPY GAP
% ===================================================================
\section{The Witness Entropy Gap}
\label{sec:gap}

\subsection{NP witnesses as compressed information}

For an NP language $L$, every yes-instance $x \in L$ has a polynomial-length witness $w$ verifiable in polynomial time:
\begin{equation}
x \in L \iff \exists w,\, |w| \leq \poly(|x|),\, V(x, w) = 1\,.
\end{equation}

The conditional complexity of the witness given the instance measures how much \textit{additional} information the witness contains beyond the instance:
\begin{equation}
K(w \mid x) = \text{``extra information needed to produce } w \text{ from } x\text{''}\,.
\end{equation}

\subsection{What P = NP would imply}

\begin{proposition}[Witness compression under P = NP]
\label{prop:compression}
If $\PP = \NP$, then for every NP language $L$ and every $x \in L$, a valid witness $w$ can be produced from $x$ by a polynomial-time algorithm $A$:
\begin{equation}
w = A(x)\,.
\end{equation}
Therefore, for some polynomial $q$ depending on the running time of $A$:
\begin{equation}
K^{q(|x|)}(w \mid x) \leq K(A) + O(\log|x|) = O(\log|x|)\,,
\label{eq:compression}
\end{equation}
since the algorithm $A$ is a fixed finite program running in polynomial time, and $O(\log|x|)$ bits encode the input length.
\end{proposition}

\begin{proof}
If $\PP = \NP$, the search version of SAT is solvable in polynomial time (by self-reducibility). Given a satisfiable formula $\varphi$, the algorithm $A$ finds a satisfying assignment bit-by-bit in polynomial time. The program ``run $A$ on input $x$'' has length $|A| + O(\log|x|)$, where $|A|$ is a constant.
\end{proof}

\subsection{The entropy obstruction}

\begin{remark}[Unbounded vs.\ resource-bounded complexity -- why $K^t$ is essential]
\label{rem:bounded-vs-unbounded}
The distinction between $K(w \mid x)$ and $K^t(w \mid x)$ is crucial and must be understood before the main theorem. For any decidable language $L$ and any $(x, w)$ with $V(x,w) = 1$, the brute-force verifier-search has $K(w \mid x) \leq O(1)$ (the search program has constant length -- it simply enumerates all candidates and checks each). The information-theoretic content of NP-hardness lies entirely in the \textit{time} required to extract the witness, not in the \textit{description length}. This is why $K^t$ with polynomial $t$ is the correct measure: it captures the computational cost of witness extraction, which is the essence of the $\PP$ vs $\NP$ question.
\end{remark}

\begin{theorem}[Existence of high-entropy witnesses -- resource-bounded, conditional]
\label{thm:high_entropy}
If $\PP \neq \NP$, then for every NP-complete language $L$, there exist infinitely many yes-instances $x \in L$ such that for every polynomial $q$, \textbf{every} valid witness $w$ for $x$ satisfies:
\begin{equation}
K^{q(|x|)}(w \mid x) \geq |w| - O(\log |x|)\,.
\label{eq:high_entropy}
\end{equation}
That is, no polynomial-time program given $x$ can produce $w$. (The ``infinitely many'' follows because finitely many exceptions can be hardcoded as a constant-size lookup table, leaving a polynomial-time algorithm for all remaining instances.)

Equivalently (by contraposition): if there exist a polynomial $q$ and a constant $c$ such that for every NP-complete language $L$ and every instance $x \in L$, some valid witness $w$ satisfies $K^{q(|x|)}(w \mid x) \leq c$, then $\PP = \NP$.

Note on quantifier order: the theorem asserts that a \emph{single} infinite set of hard instances resists all polynomials $q$ simultaneously. This follows because $K^{q_1}(w \mid x) \leq K^{q_2}(w \mid x)$ whenever $q_1 \geq q_2$ (more time can only decrease description length). Therefore, if an instance $x$ has a witness $w$ with $K^{q_0(|x|)}(w \mid x) \geq |w| - O(\log |x|)$ for some polynomial $q_0$, the same bound holds for all $q \leq q_0$. The contrapositive guarantees that for every fixed $q$, hard instances exist; the monotonicity of $K^t$ in $t$ ensures that instances that are hard for large time bounds are automatically hard for all smaller ones, so the infinite set can be chosen independently of $q$.
\end{theorem}

\begin{proof}
We prove the contrapositive. Suppose for some NP-complete language $L$, some polynomial $q$, and some constant $c$, every yes-instance $x \in L$ has a witness $w$ with $K^{q(|x|)}(w \mid x) \leq c$. Then for every yes-instance $x \in L$, there exists a program $p_x$ of length $\leq c$ that, given $x$, produces a valid witness $w$ in time $q(|x|)$. (Note: $p_x$ may depend on $x$, but there are at most $2^{c+1} - 1$ programs of length $\leq c$.) This gives a polynomial-time search algorithm for $L$: on input $x$, enumerate all programs of length $\leq c$, run each with input $x$ for at most $q(|x|)$ steps, and verify each output. Since $L$ is NP-complete, every NP search problem reduces to the search version of $L$ via polynomial-time reductions. A polynomial-time search algorithm for $L$ therefore yields polynomial-time search algorithms for all NP search problems, implying $\mathrm{FP} = \mathrm{FNP}$ (where $\mathrm{FP}$ and $\mathrm{FNP}$ are the function-problem analogues of $\PP$ and $\NP$) and hence $\PP = \NP$.

\textit{Existence of infinitely many high-entropy instances (given $\PP \neq \NP$):} The formal proof is the contrapositive above. The \textit{existence} of high-entropy instances is guaranteed non-constructively: if every SAT instance had a low-entropy witness ($K^{q(n)}(w \mid x) \leq c$ for some fixed $q, c$), we would have $\PP = \NP$. Under $\PP \neq \NP$, such instances must exist. Moreover, there must be \textit{infinitely many}: if only finitely many hard instances existed, they could be hardcoded into a constant-size lookup table, leaving a polynomial-time algorithm for all remaining instances -- contradicting $\PP \neq \NP$. Constructing the hard instances explicitly is non-trivial -- a naive approach via Cook-Levin encoding of ``verify $v = r$'' for a random string $r$ fails because $r$ is recoverable from the formula structure in polynomial time, giving $K^{\poly}(r \mid \varphi_r) = O(1)$. The extension to any NP-complete language $L$ follows from NP-completeness: since SAT reduces to $L$ in polynomial time, a polynomial-time witness-finding algorithm for $L$ composed with the reduction yields a polynomial-time witness-finding algorithm for SAT (one reduces a SAT instance to an $L$-instance, finds a witness for $L$, and uses the witness to reconstruct a satisfying assignment for the original SAT instance).

\textit{Why unbounded $K$ fails:} For any SAT instance $x$ with a satisfying assignment $w$, the brute-force program ``enumerate all assignments, test each, output the first satisfying one'' has length $O(1)$ and finds $w$ from $x$ in time $O(2^{|x|})$. Thus $K(w \mid x) \leq O(1)$ -- the unbounded conditional complexity is trivially small. This is why the resource bound is essential: the information-theoretic content of NP-hardness lies in the \textit{time} required to extract the witness, not in the description length.
\end{proof}


\subsection{The gap}

Combining Proposition~\ref{prop:compression} and Theorem~\ref{thm:high_entropy}:

\begin{corollary}[Witness Entropy Gap -- the core contradiction structure]
\label{cor:gap}
Consider the infinitely many hard instances whose existence is guaranteed by Theorem~\ref{thm:high_entropy} (under $\PP \neq \NP$). For unique-witness instances, the contradiction structure is sharpest:
\begin{itemize}
\item $\PP = \NP$ $\Rightarrow$ $K^{q(|x|)}(w \mid x) \leq O(\log |x|)$ for some polynomial $q$ (Proposition~\ref{prop:compression}), since any algorithm $A$ must output the unique witness $w = A(x)$.
\item $\PP \neq \NP$ $\Rightarrow$ $K^{q(|x|)}(w \mid x) \geq |w| - O(\log |x|)$ for all polynomials $q$ (Theorem~\ref{thm:high_entropy}).
\end{itemize}
This is not a proof of $\PP \neq \NP$, but a sharp characterization: $\PP = \NP$ if and only if all NP-witnesses are polynomially compressible. The entropy gap $\Omega(|w|)$ vs.\ $O(\log n)$ quantifies exactly what separating $\PP$ from $\NP$ requires.
\end{corollary}

\begin{remark}[Unique vs.\ multiple witnesses]
Theorem~\ref{thm:high_entropy} guarantees instances where \textit{every} valid witness has high conditional complexity. For these instances, the contradiction structure holds regardless of how many witnesses exist: any algorithm $A$ must output \textit{some} valid witness, and \textit{all} of them satisfy the high-entropy bound~\eqref{eq:high_entropy}.

For unique-witness instances, the structure is especially clean: $A$ has no choice over which witness to produce. In general, however, the high-entropy instances guaranteed by Theorem~\ref{thm:high_entropy} exist non-constructively (via contraposition), and it remains open whether they can be constructed explicitly -- see Section~\ref{sec:obstruction}.
\end{remark}


% ===================================================================
% 5. ALGORITHMIC POSITIVITY
% ===================================================================
\section{Algorithmic Positivity}
\label{sec:positivity}

We now formalize the No-Go framework.

\subsection{Definition}

\begin{definition}[Algorithmic Positivity]
\label{def:alg_pos}
An NP language $L$ (with verifier $V$) is \textit{algorithmically positive} if there exists a \textit{witness-producing} function $W: \{0,1\}^* \to \{0,1\}^*$ satisfying (AP1) and (AP2):

\begin{enumerate}
\item[\textbf{(AP1)}] \textbf{Polynomial scaling.} $W$ is computable in time $\poly(|x|)$, and for every $x \in L$: $V(x, W(x)) = 1$.

\item[\textbf{(AP2)}] \textbf{Entropic convexity.} There exists a polynomial $q$ such that for every $x \in L$:
\begin{equation}
K^{q(|x|)}(W(x) \mid x) \leq O(\log|x|)\,.
\end{equation}
The witness production is ``entropically cheap'': the witness $W(x)$ adds at most logarithmic resource-bounded information beyond the input. (Note: since $W$ is a fixed polynomial-time program, the polynomial $q$ is determined by $W$ and does not depend on the instance $x$.)
\end{enumerate}
\end{definition}

\begin{remark}
Note that (AP2) targets the \textit{witness} $W(x)$, not just the decision bit. A single decision bit $D(x) \in \{0,1\}$ trivially satisfies $K(D(x) \mid x) \leq 1$, so entropic convexity would be vacuous for decision functions. The meaningful measure is the resource-bounded complexity of producing a \textit{witness} -- this is where the information-theoretic content of NP-hardness resides.
\end{remark}

\begin{remark}
Every $\PP$-decidable NP language is algorithmically positive: if $L \in \PP \cap \NP$, then the polynomial-time search algorithm (obtained via self-reducibility for NP-complete problems, or directly for natural NP problems) produces a witness $W(x)$ with $K^{q(|x|)}(W(x) \mid x) \leq |W_{\text{prog}}| + O(\log|x|) = O(\log|x|)$.
\end{remark}

\begin{remark}[Heuristic properties of P-decidable NP languages]
\label{rem:heuristic-ap}
For NP languages in $\PP$, the witness-producing function $W$ additionally satisfies two heuristic properties that capture the ``structural simplicity'' of polynomial-time witness extraction:

\textbf{(H1) Compositional stability.} For structured instances under a well-defined composition operator $\circ$, the witness for the composed input is typically determined by the component witnesses with at most logarithmic overhead: $K^{\poly}(W(x_1 \circ x_2) \mid W(x_1), W(x_2)) \leq O(\log(|x_1|+|x_2|))$.

\textbf{(H2) Monotone stability.} Small perturbations of the input produce witnesses of bounded additional complexity: $d_H(x,x') \leq 1 \implies K(W(x) \mid W(x'), x) \leq O(\mathrm{poly}(\log|x|))$.

These properties are motivated by the structure of $\NP$-complete problems but are not used in the formal No-Go argument (Theorem~\ref{thm:nogo}). They serve as additional heuristic evidence for the Entropic Separation Conjecture.
\end{remark}

\subsection{The Entropic Separation Conjecture}

\begin{conjecture}[Entropic Separation -- ESC]
\label{conj:entropic}
No NP-complete language is algorithmically positive. That is: for every NP-complete language $L$ and every polynomial-time witness-producing function $W$, there exist infinitely many yes-instances $x \in L$ such that $W$ fails to produce a valid witness for $x$ (i.e., $V(x, W(x)) = 0$), or $K^{q(|x|)}(W(x) \mid x) \geq \omega(\log|x|)$ for all polynomials $q$.

An equivalent formulation in terms of truth-table complexity (see Theorem~\ref{thm:slice_nogo}): for some NP-complete language $L$, for every polynomial $q$, infinitely many $n$ satisfy $K^{q(2^n)}(L_n) \geq n$. (Both formulations are equivalent to $\PP \neq \NP$; their mutual equivalence follows transitively.)
\end{conjecture}

\begin{proposition}
\label{prop:equiv}
The Entropic Separation Conjecture~\ref{conj:entropic} implies $\PP \neq \NP$.
\end{proposition}

\begin{proof}
If $\PP = \NP$, then for every NP-complete language $L$, the search version is solvable in polynomial time (for SAT, by self-reducibility; for general NP-complete $L$, by reducing $L$ to SAT, applying the SAT search algorithm, and reconstructing the witness). This gives a polynomial-time witness-producing function $W$ satisfying (AP1). Since $W$ is a fixed program running in polynomial time, $K^{q(|x|)}(W(x) \mid x) \leq |W_{\text{prog}}| + O(\log|x|) = O(\log|x|)$, satisfying (AP2). Therefore, NP-complete languages would be algorithmically positive, contradicting the conjecture.
\end{proof}

\begin{remark}[ESC is equivalent to $\PP \neq \NP$]
\label{rem:esc-equivalence}
The converse of Proposition~\ref{prop:equiv} also holds: $\PP \neq \NP$ implies ESC. Indeed, if $\PP = \NP$, then every NP-complete language $L$ has a polynomial-time algorithm, giving $K^{\mathrm{poly}(2^n)}(L_n) = O(\log n)$ by Lemma~\ref{lem:p_low_entropy}, which violates the ESC condition $K^{q(2^n)}(L_n) \geq n$. Thus ESC $\Leftrightarrow$ $\PP \neq \NP$: the Entropic Separation Conjecture is a \textit{reformulation} of the P vs NP problem in information-theoretic language, not an independent strengthening.
\end{remark}

\subsection{Why NP-complete problems resist algorithmic positivity}

The failure of each condition for NP-complete problems:

\textbf{(AP2) fails:} For NP-complete problems, producing a witness $w$ for $x \in L$ requires extracting information that is algorithmically random relative to the instance. The resource-bounded conditional complexity $K^t(w \mid x)$ can be $\Omega(|w|)$ (Theorem~\ref{thm:high_entropy}, conditional on $\PP \neq \NP$). No polynomial-time witness-producing function $W$ can achieve $K^{\poly}(W(x) \mid x) \leq O(\log|x|)$ for all instances.

As additional heuristic evidence (not part of the formal argument), the properties (H1)--(H2) from Remark~\ref{rem:heuristic-ap} also fail for NP-complete problems:

\textbf{(H1) fails:} SAT is not compositionally decomposable in the required sense. The satisfiability of a conjunction $\varphi_1 \wedge \varphi_2$ is not determined by the satisfiability of $\varphi_1$ and $\varphi_2$ individually (they share variables). This non-decomposability is precisely the source of NP-hardness.

\textbf{(H2) fails:} SAT exhibits extreme sensitivity in its search version: flipping a single clause can restructure the entire witness space. Near the phase transition, a single bit-flip can fragment the solution cluster into exponentially many components, producing witnesses of unbounded additional complexity.


% ===================================================================
% 6. THE NO-GO THEOREM
% ===================================================================
\section{The No-Go Theorem}
\label{sec:nogo}

\begin{theorem}[Entropic No-Go -- Reformulation]
\label{thm:nogo}
The following are equivalent:
\begin{enumerate}
\item[(i)] $\PP \neq \NP$.
\item[(ii)] The Entropic Separation Conjecture~\ref{conj:entropic} holds.
\item[(iii)] No NP-complete language is decidable in polynomial time. That is: for every NP-complete language $L$ and every polynomial-time algorithm $A$, there exist instances $x$ where $A(x) \neq L(x)$.
\end{enumerate}
The equivalence $(i) \Leftrightarrow (ii)$ shows that $\PP \neq \NP$ is \textit{equivalent} to an entropy obstruction: polynomial-time algorithms cannot compress NP-witnesses below their intrinsic Kolmogorov complexity.
\end{theorem}

\begin{proof}
$(ii) \Rightarrow (i)$: Shown in Proposition~\ref{prop:equiv}. Contrapositive: if $\PP = \NP$, then every NP-complete language has a polynomial-time witness-producing function (via the SAT search algorithm and polynomial-time reductions), satisfying both (AP1) and (AP2). This makes NP-complete languages algorithmically positive, contradicting ESC.

$(i) \Rightarrow (ii)$: Contrapositive: suppose ESC fails, i.e., some NP-complete language $L$ is algorithmically positive. Then $L$ has a polynomial-time witness-producing function $W$ with $K^{\poly}(W(x) \mid x) \leq O(\log|x|)$. In particular, $W$ is a polynomial-time search algorithm for $L$. Since $L$ is NP-complete, this implies $\PP = \NP$.

$(i) \Leftrightarrow (iii)$: (iii) states that for every NP-complete language $L$, $L \notin \PP$. This is equivalent to $\PP \neq \NP$ by the existence of NP-complete languages (if any NP-complete $L$ were in $\PP$, then $\PP = \NP$; conversely, $\PP = \NP$ would place every NP-complete language in $\PP$).

The content of the theorem is not a new separation result, but the \textit{reformulation}: $\PP \neq \NP$ is equivalent to an entropic incompressibility condition on NP-witnesses. This reformulation identifies the precise structure (witness entropy) and the precise gap (uniformity bridge) that any proof must address.
\end{proof}

\subsection{Exclusion of alternative scenarios}

\begin{proposition}[Excluded scenarios]
\label{prop:excluded}
The following scenarios are incompatible with the conjunction of entropic positivity and Kolmogorov obstruction:

\begin{enumerate}
\item[\textbf{(X1)}] \textbf{Polynomial algorithm with exponential entropy production.} An algorithm that ``invents'' new information during computation. Excluded because for any deterministic algorithm $A$ running in time $t_A(|x|)$, the resource-bounded conditional complexity satisfies $K^{t_A(|x|)}(A(x) \mid x) \leq |A| + O(\log|x|)$: the program ``run $A$ on input $x$'' has constant description length. In particular, the \textit{conditional} algorithmic information of $A(x)$ given $x$ is at most $O(1)$ bits at the time scale of $A$'s own computation -- meaning $A$ cannot ``create'' information that was not already implicit in the input. (For randomized algorithms, the random bits $r$ provide additional input, so $K(A(x,r) \mid x)$ can be high. However, the P vs NP question concerns deterministic computation; moreover, under widely believed derandomization assumptions, $\BPP = \PP$~\cite{ImpagliazzoWigderson1997}, so randomness does not provide additional power.)

\item[\textbf{(X2)}] \textbf{Compositional collapse of NP.} All SAT instances decompose into independently solvable parts. Excluded by the NP-completeness of SAT: if SAT decomposed, then all NP problems would decompose (via reduction), but UNIQUE-SAT instances with random solutions do not decompose. (The relevance of UNIQUE-SAT is supported by the Valiant-Vazirani theorem~\cite{ValiantVazirani1986}: SAT reduces to UNIQUE-SAT under randomized reductions.)

\item[\textbf{(X3)}] \textbf{Smooth decision landscape.} All SAT instances have monotonically stable decision boundaries. Excluded by phase-transition phenomena in random SAT: near the satisfiability threshold, the decision changes discontinuously under small clause additions.
\end{enumerate}
\end{proposition}


% ===================================================================
% 7. RESOURCE-BOUNDED SLICE ENTROPY
% ===================================================================
\section{Resource-Bounded Slice Entropy}
\label{sec:slices}

The witness entropy gap of Section~\ref{sec:gap} targets individual instances. A cleaner formulation targets the \textit{truth table} of the language on all inputs of a given length.

\subsection{Slices of a language}

\begin{definition}[Slice]
For a language $L \subseteq \{0,1\}^*$, the $n$-th \textit{slice} $L_n \in \{0,1\}^{2^n}$ is the truth table of $L$ on inputs of length $n$: $(L_n)_i = 1 \iff x_i \in L$, where $x_1, \ldots, x_{2^n}$ are the $n$-bit strings in lexicographic order.
\end{definition}

\begin{definition}[Resource-bounded algorithmic entropy]
For a time bound $t(\cdot)$:
\begin{equation}
K^t(x) := \min\{|p| : U(p) = x \text{ in at most } t(|x|) \text{ steps}\}\,.
\label{eq:Kt}
\end{equation}
\end{definition}

\subsection{The slice entropy lemma}

\begin{lemma}[P implies low slice entropy]
\label{lem:p_low_entropy}
If $L \in \PP$, then there exists a polynomial $q$ and a constant $c$ such that for all $n$:
\begin{equation}
K^{q(2^n)}(L_n) \leq c + \lceil \log n \rceil\,.
\label{eq:p_entropy}
\end{equation}
\end{lemma}

\begin{proof}
Let $A$ be a polynomial-time algorithm deciding $L$ in time $n^d$. The program $p$: ``read $n$ from input; for each $x \in \{0,1\}^n$, run $A(x)$ and output the result bit'' has length $|A| + O(\log n)$ (to encode $n$). Its runtime is $2^n \cdot n^d = \poly(2^n)$. Thus $K^{q(2^n)}(L_n) \leq |A| + O(\log n) = c + \lceil \log n \rceil$.
\end{proof}

\subsection{The No-Go theorem via slice entropy}

\begin{theorem}[Slice Entropy No-Go]
\label{thm:slice_nogo}
If there exists an $\NP$-complete language $L$ satisfying the \textit{Entropy Hardness Hypothesis}:

\textbf{(EH)} For every polynomial $q$, there exist infinitely many $n$ with
\begin{equation}
K^{q(2^n)}(L_n) \geq n\,,
\label{eq:eh}
\end{equation}

then $\PP \neq \NP$.
\end{theorem}

\begin{proof}
Immediate from Lemma~\ref{lem:p_low_entropy}: $L \in \PP$ would give $K^{\poly(2^n)}(L_n) = O(\log n)$, contradicting~\eqref{eq:eh} for large $n$.
\end{proof}

\begin{remark}
The Entropy Hardness Hypothesis (EH) is a concrete, well-defined conjecture about a specific NP-complete language (e.g., SAT). It asserts that the truth table of SAT on $n$-bit formulas cannot be compressed to fewer than $n$ bits, even when the compression algorithm is allowed $\poly(2^n)$ time.

This is the \textit{cleanest} No-Go formulation: the ``easy direction'' ($L \in \PP \Rightarrow$ low entropy) is a theorem; the ``hard direction'' (EH for SAT) is the entire content of $\PP \neq \NP$.

Note that the threshold $n$ in EH is much smaller than the truth table length $2^n$: a random string of length $2^n$ has $K \approx 2^n$. The threshold $n$ is chosen because Lemma~\ref{lem:p_low_entropy} gives $K^{\poly(2^n)}(L_n) = O(\log n)$ for $L \in \PP$; any threshold $\omega(\log n)$ would suffice to separate from P, and $n$ provides a clean, moderate lower bound that is already superpolynomial in the input length.
\end{remark}

\begin{remark}[Encoding sensitivity of EH]
\label{rem:encoding}
The Entropy Hardness Hypothesis depends on the encoding of SAT instances as binary strings. Different encodings yield truth tables $\SAT_n$ of different lengths and potentially different Kolmogorov complexities. The hypothesis should be understood as: there exists a \emph{standard} polynomial-time encoding $\mathrm{enc}: \mathrm{Formulas} \to \{0,1\}^*$ such that $K^{q(2^n)}(\SAT_n^{\mathrm{enc}}) \geq n$ for infinitely many $n$. For resource-bounded $K^t$, changing the universal machine introduces a polynomial slowdown rather than just an additive constant. However, since EH quantifies over \emph{all} polynomials $q$, a polynomial slowdown in the time bound does not affect the hypothesis: if $K^{q(2^n)}(\SAT_n) \geq n$ for all $q$, the same holds after re-parameterization. The \emph{length} of the truth table (and hence the threshold $n$) is encoding-dependent.
\end{remark}

\begin{remark}[Connection to circuit lower bounds]
\label{rem:circuit-connection}
The Entropy Hardness Hypothesis is closely related to circuit lower bounds. A Boolean function $f_n$ computed by a circuit of size $s$ has truth table describable in $O(s \log s)$ bits, so $K^{O(s \cdot 2^n)}(f_n) \leq O(s \log s) + O(\log n)$. Thus:
\begin{itemize}
\item EH with threshold $n$ implies $\SAT_n$ requires circuits of size $s \geq 2^{\Omega(n / \log n)}$, i.e., $\NP \not\subseteq \mathrm{SIZE}(2^{o(n)})$.
\item In particular, EH implies $\NP \not\subseteq \PP/\poly$ (since polynomial-size circuits give $K^{\poly(2^n)}(f_n) \leq O(n^c \log n)$, which is $o(n)$ for large $n$).
\end{itemize}
This shows that the entropy framework and the circuit complexity approach are two perspectives on the same underlying separation.
\end{remark}

\subsection{Why (EH) should be true}

The truth table $\SAT_n$ has length $2^n$ and describes the satisfiability of all $n$-bit encodings of Boolean formulas. There are $2^{2^n}$ possible truth tables of this length, and $\SAT_n$ is a \textit{specific} one determined by a complete mathematical definition ($\SAT$). Nevertheless:

\begin{enumerate}
\item \textbf{Pseudorandom structure.} The satisfiability boundary in the space of formulas exhibits phase-transition behavior \cite{Monasson1999}: near the threshold, the pattern of satisfiable vs.\ unsatisfiable instances has high apparent randomness.

\item \textbf{No known compression.} Despite decades of SAT-solver research, no method compresses $\SAT_n$ to $o(2^n)$ bits in $\poly(2^n)$ time. The best algorithms (CDCL solvers) run in time $2^{\Theta(n)}$ in the worst case.

\item \textbf{Cryptographic evidence.} If (EH) fails, then efficient PRGs from one-way functions would not exist \cite{ImpagliazzoWigderson1997}, collapsing much of modern cryptography.
\end{enumerate}


% ===================================================================
% 8. THE PRECISE OBSTRUCTION
% ===================================================================
\section{The Precise Obstruction}
\label{sec:obstruction}

\subsection{The uniformity gap}

The main argument (Section~\ref{sec:gap}) has a precise gap that we now identify.

\textbf{What is proven:} There \textit{exist} NP instances whose witnesses have high conditional Kolmogorov complexity (Theorem~\ref{thm:high_entropy}). This is a \textit{non-uniform} statement about individual strings.

\textbf{What is needed:} A \textit{uniform} statement: for every polynomial-time algorithm $A$, there exist infinitely many instances $x$ where $A(x)$ fails. This requires converting the existential entropy bound into an \textit{adversarial} bound against all efficient algorithms simultaneously.

\begin{definition}[The uniformity bridge]
\label{def:uniformity}
Let $L$ be an NP-complete language with verifier $V$. The \textit{uniformity bridge} is the statement: for every polynomial-time computable function $A: \{0,1\}^* \to \{0,1\}^*$, there exist infinitely many yes-instances $x \in L$ such that $A$ fails to produce a valid witness:
\begin{equation}
\forall A \in \mathrm{FP},\, \exists^\infty x \in L:\, V(x, A(x)) = 0\,.
\end{equation}
Equivalently, in entropic terms: for every polynomial-time $A$ and every polynomial $q$, there exist infinitely many $x \in L$ such that every valid witness $w$ for $x$ satisfies $K^{q(|x|)}(w \mid x) > |A| + O(\log|x|)$.

Note: the uniformity bridge is formulated as a search statement ($A$ must produce a \textit{witness}, i.e., $A \in \mathrm{FP}$ and $V(x, A(x)) = 1$). The connection to the decision problem $\PP \neq \NP$ follows from the NP-completeness of $L$: for NP-complete languages, the search and decision versions are polynomial-time equivalent (via self-reducibility for SAT, or more generally via Cook-Levin plus polynomial-time reductions).
\end{definition}

The uniformity bridge is the P-vs-NP analog of:
\begin{itemize}
\item The ``higher Gross-Zagier formula'' in BSD (coupling analytic rank to algebraic rank).
\item The ``hard direction'' in the Hodge conjecture (showing positivity of the Hodge-Riemann form implies algebraicity).
\end{itemize}

In each case, the positivity structure is established, but a \textit{coupling theorem} is needed to convert positivity into a definitive result.

\subsection{Known partial results}

The uniformity bridge is partially established in restricted settings:

\begin{enumerate}
\item \textbf{Resolution proof systems:} Exponential lower bounds on resolution refutations of random $k$-SAT instances near the threshold \cite{BenSasson2001}. This proves that resolution-based SAT solvers cannot solve these instances efficiently.

\item \textbf{Bounded-depth circuits:} $\SAT \notin \mathrm{AC}^0$ (Furst-Saxe-Sipser 1984, Ajtai 1983; see~\cite{Hastad1987} for optimal bounds). This proves the entropy obstruction for constant-depth circuits.

\item \textbf{Monotone circuits:} Exponential lower bounds for monotone circuits computing CLIQUE \cite{Razborov1985}; related bounded-depth lower bounds \cite{Razborov1987}.

\item \textbf{Algebraic proof systems:} Exponential lower bounds for Nullstellensatz and Polynomial Calculus refutations \cite{Grigoriev2001}.
\end{enumerate}

Each of these is a \textit{model-specific} instantiation of the entropy obstruction. The challenge is to prove it for \textit{general} polynomial-time computation -- which is precisely the $\PP \neq \NP$ problem.

\subsection{Why the gap is hard}

The uniformity gap is hard for the same reason that $\PP$ vs $\NP$ is hard: it requires a statement about \textit{all} polynomial-time algorithms. In the uniform model, this is a universal quantifier over a countably infinite, recursively enumerable set (all Turing machines running in polynomial time). In the non-uniform model ($\PP/\poly$), the quantifier ranges over all polynomial-size circuit families, an uncountable set.

The entropy approach \textit{reduces} the problem but does not \textit{solve} it:
\begin{quote}
\textit{P vs NP is equivalent to asking whether Kolmogorov complexity can be polynomially compressed for NP-witness structures. The entropy framework identifies what must be shown (witness incompressibility) and why it should be true (random witnesses exist), but converting this to a proof against all algorithms remains the central challenge.}
\end{quote}

\subsection{Four candidate bridge strategies}

We identify four concrete research directions for closing the uniformity gap:

\textbf{(B1) Dense families with unique random witnesses.} Construct \textit{dense} families $S_n \subseteq \{0,1\}^{\poly(n)}$ with $|S_n|/2^{\poly(n)}$ non-vanishing, such that every $x \in S_n$ has a unique valid witness $w_x$ with $K^{\poly(n)}(w_x \mid x) \geq |w_x| - O(\log n)$. Density ensures that no polynomial-time algorithm can ``avoid'' all hard instances. The construction requires efficient reductions from Kolmogorov-random strings to SAT instances with unique solutions -- a non-trivial but plausible coding-theoretic task.

\textbf{(B2) Self-referential diagonalization.} For each polynomial-time machine $A$, construct an instance $x_A$ whose unique valid witness $w$ contains (or encodes) a description of $A$ itself, so that $K^{\poly(|x_A|)}(w \mid x_A) \geq |A|$. This creates a self-referential obstruction: $A$ cannot produce a witness that is ``about itself'' without exceeding its own description length. The approach combines G\"odel-style diagonalization with Kolmogorov complexity and avoids relativization because the self-reference targets the machine's \textit{description}, not its oracle behavior.

\textbf{(B3) Resource-bounded Kolmogorov lower bounds on slices.} Prove that for SAT (or another NP-complete language), the resource-bounded complexity satisfies $K^{\poly(2^n)}(\SAT_n) \geq 2^{n-o(n)}$ for infinitely many $n$. This is a \textit{uniform} statement about the truth table and directly implies $\PP \neq \NP$ via the Slice Entropy No-Go (Theorem~\ref{thm:slice_nogo}). The strategy connects to proof complexity: if every short proof of an instance's satisfiability has high Kolmogorov complexity, then no polynomial-time algorithm can systematically produce such proofs.

\textbf{(B4) The Minimum Circuit Size Problem (MCSP).} Recent work connects resource-bounded Kolmogorov complexity to circuit complexity via the \textit{Minimum Circuit Size Problem} ($\mathrm{MCSP}$): given a truth table $t \in \{0,1\}^{2^n}$ and a parameter $s$, decide whether there exists a circuit of size $\leq s$ computing $t$. Allender et al.~\cite{Allender2006} established deep connections between $\mathrm{MCSP}$ and $K^t$. While $\mathrm{MCSP} \in \NP$, its complexity status is open: it is not known to be $\NP$-complete, yet proving $\mathrm{MCSP} \in \PP$ would have strong derandomisation consequences. If $\mathrm{MCSP}$ is hard for some complexity class (e.g., $\mathrm{MCSP} \notin \PP/\mathrm{poly}$), this would yield circuit lower bounds that imply the Entropy Hardness Hypothesis. The chain is indirect but concrete: \textit{hardness of MCSP $\Rightarrow$ circuit lower bounds $\Rightarrow$ truth-table incompressibility $\Rightarrow$ EH}.


% ===================================================================
% 9. CONNECTION TO PROOF COMPLEXITY
% ===================================================================
\section{Proof Complexity as a Testing Ground}
\label{sec:proof_complexity}

\subsection{SAT algorithms and short proofs}

A deep connection exists between algorithms and proofs:
\begin{equation}
\SAT \in \PP \implies \text{every unsatisfiable formula has a short refutation}\,,
\end{equation}
in a proof system at least as strong as Extended Frege.

\begin{proposition}[Proof complexity formulation -- Cook {\cite{Cook1975}}]
\label{prop:proof_complexity}
If $\PP \neq \NP$, then there is no polynomially bounded proof system for tautologies. Equivalently: $\PP = \NP$ if and only if there exists a proof system in which every tautology of length $n$ has a proof of length $\poly(n)$.
\end{proposition}

This is Cook's reformulation~\cite{Cook1975}: proving superpolynomial lower bounds on proof length in every propositional proof system is equivalent to showing $\PP \neq \coNP$ (which follows from $\PP \neq \NP$).

\subsection{Entropy of proofs}

The entropy framework connects to proof complexity via the following question: for a random unsatisfiable formula $\varphi$ and its shortest refutation $\pi$ in a given proof system, what is $K(\pi \mid \varphi)$? If $\pi$ is ``algorithmically random'' relative to $\varphi$ -- meaning $K(\pi \mid \varphi)$ is close to $|\pi|$ -- then no efficient algorithm can systematically produce such refutations.

This provides a concrete testing ground: if one can show that for random unsatisfiable $3$-SAT near the threshold, the shortest refutation in Extended Frege has both superpolynomial length and high conditional Kolmogorov complexity, this would be strong evidence for the entropy obstruction.

\subsection{Entropy proof of $\SAT \notin \mathrm{AC}^0$}
\label{sec:ac0}

As a concrete instantiation of the entropy framework, we give an entropy-based argument for the classical result $\mathrm{PAR} \notin \mathrm{AC}^0$ (Furst-Saxe-Sipser 1984, Ajtai 1983). Since $\mathrm{AC}^0$ is closed under $\mathrm{AC}^0$-reductions and parity is $\mathrm{AC}^0$-reducible to SAT (via a standard reduction that encodes the parity constraint as a CNF formula using constant-depth circuitry), this also implies $\SAT \notin \mathrm{AC}^0$.

\begin{proposition}[Entropy obstruction for $\mathrm{AC}^0$]
\label{prop:ac0}
Let $\{C_n\}$ be a constant-depth circuit family of polynomial size $s(n) = n^{O(1)}$ computing a Boolean function $f_n: \{0,1\}^n \to \{0,1\}$. Then the truth table $f_n \in \{0,1\}^{2^n}$ satisfies
\begin{equation}
K^{s(n) \cdot 2^n}(f_n) \leq O(n^c \log n)
\end{equation}
for some constant $c$ depending on the depth and size bound. In contrast, a \textit{random} Boolean function $g_n$ satisfies $K(g_n) \geq 2^n - O(1)$ with high probability. The entropy gap between $\mathrm{AC}^0$-computable functions and random functions is thus exponential: $O(n^c \log n)$ vs.\ $2^n$.
\end{proposition}

\begin{proof}[Proof sketch]
A depth-$d$ circuit of size $s$ on $n$ inputs is described by $O(s \cdot \log s)$ bits (gate types and wiring). Given this description, the truth table on all $2^n$ inputs is computed in time $O(s \cdot 2^n)$ by brute-force evaluation. Thus $K^{O(s \cdot 2^n)}(f_n) \leq O(s \cdot \log s) + O(\log n)$. For polynomial size $s = n^c$: $K^{n^c \cdot 2^n}(f_n) \leq O(n^c \log n)$.

The classical result $\mathrm{PAR}_n \notin \mathrm{AC}^0$ (Furst-Saxe-Sipser 1984, Ajtai 1983, H\r{a}stad 1987) shows that the parity function cannot be computed by polynomial-size constant-depth circuits. In the entropy framework, this means: any circuit computing $\mathrm{PAR}_n$ requires super-polynomial size $s(n) = 2^{n^{\Omega(1/d)}}$, so $K^{\mathrm{poly}(2^n)}(\mathrm{PAR}_n)$ is \textit{not} bounded by $O(n^c \log n)$ when computed via $\mathrm{AC}^0$ circuits. The entropy method does not \textit{reprove} $\mathrm{PAR} \notin \mathrm{AC}^0$, but it \textit{reframes} the separation: $\mathrm{AC}^0$-computable truth tables occupy a negligible fraction of the space of all Boolean functions, and parity lies outside this low-entropy region.
\end{proof}

This demonstrates the entropy perspective on a known separation. For general $\PP/\mathrm{poly}$, the analogous bound would give $K^{\mathrm{poly}(2^n)}(f_n) \leq O(n^c \log n)$ for any $f \in \PP/\mathrm{poly}$. Proving that $\SAT_n$ exceeds this bound would establish $\NP \not\subseteq \PP/\mathrm{poly}$ -- a major open problem that implies $\PP \neq \NP$.

\begin{remark}[H\r{a}stad's bound as entropy capacity]
\label{rem:hastad-entropy}
H\r{a}stad's switching lemma~\cite{Hastad1987} can be interpreted as an \emph{information-theoretic capacity bound}: after applying a random restriction that fixes each variable independently with probability $1-p$, a depth-$d$ $\mathrm{AC}^0$ circuit on the surviving variables simplifies to a decision tree of depth $O(\log n)^{d-1}$. In entropy terms: each layer of an $\mathrm{AC}^0$ circuit can ``aggregate'' at most $O(\log n)$ bits of information from its inputs (via the fan-in bound). After $d$ layers, the total information capacity is $O((\log n)^d)$ bits -- far below the $\Omega(n)$ bits needed to compute parity. This capacity bound $O((\log n)^d) \ll n$ is the entropy-theoretic core of the $\mathrm{PAR} \notin \mathrm{AC}^0$ separation.
\end{remark}

\subsection{Geometric Complexity Theory}

Mulmuley's Geometric Complexity Theory (GCT) \cite{Mulmuley2011} approaches $\PP$ vs $\NP$ through algebraic geometry and representation theory. The GCT program seeks \textit{obstructions} -- representation-theoretic invariants that distinguish easy from hard functions.

In the entropy framework, GCT obstructions are a specific type of \textit{algebraic entropy certificate}: they certify that a function has high ``algebraic complexity'' by exhibiting an invariant that polynomial-size circuits cannot match.

The connection to algorithmic positivity: GCT obstructions are \textit{non-natural} (they are not efficiently computable properties of truth tables) and \textit{non-relativizing} (they depend on algebraic structure). They thus satisfy specifications (S1)--(S2). Whether they satisfy (S3) (non-algebrization) is an active research question.


% ===================================================================
% 10. RELATED WORK
% ===================================================================
\section{Related Work}
\label{sec:related}

The connection between Kolmogorov complexity and computational complexity has been studied extensively. We briefly situate our work within this literature.

\textbf{Resource-bounded Kolmogorov complexity and circuit lower bounds.}
Allender et al.~\cite{Allender2006} established deep connections between the Minimum Circuit Size Problem (MCSP), resource-bounded Kolmogorov complexity $K^t$, and derandomization. They showed that $K^t$ can encode significant information about circuit complexity: a truth table $f_n$ computed by circuits of size $s$ satisfies $K^{O(s \cdot 2^n)}(f_n) \leq O(s \log s)$ (the time bound $O(s \cdot 2^n)$ accounts for evaluating the circuit on all $2^n$ inputs), linking slice entropy directly to circuit lower bounds. Allender~\cite{Allender2001} surveyed how Kolmogorov complexity and derandomization interact, arguing that understanding the complexity of ``random strings'' (strings with high $K^t$) is central to resolving P vs NP.

\textbf{MCSP and NP-hardness.}
Hirahara~\cite{Hirahara2018} showed that the NP-hardness of MCSP under certain reductions would imply new circuit lower bounds, providing a concrete path from $K^t$-based hardness to P vs NP separations. The chain ``hardness of MCSP $\Rightarrow$ circuit lower bounds $\Rightarrow$ truth-table incompressibility $\Rightarrow$ EH'' motivates our bridge strategy B4.

\textbf{Natural proofs and pseudorandomness.}
The natural proofs barrier~\cite{RazborovRudich1997} is closely connected to the non-computability of $K$: a ``natural'' property must be efficiently computable, but high Kolmogorov complexity is not. Our approach exploits this connection systematically. Impagliazzo and Wigderson~\cite{ImpagliazzoWigderson1997} showed that strong circuit lower bounds imply derandomization ($\mathrm{P} = \mathrm{BPP}$), which in Impagliazzo's ``Five Worlds'' framework~\cite{Impagliazzo1995} corresponds to Heuristica or beyond. The failure of EH would place us in Algorithmica (where P = NP), which is widely believed to be false.

\textbf{Distinction from prior work.}
The main novelty of this paper is not the use of Kolmogorov complexity in complexity theory (which is well-established), but the specific framing as a No-Go theorem via algorithmic positivity, the identification of the uniformity bridge as the precise remaining gap, and the connection to the broader positivity pattern across mathematics.

% ===================================================================
% 11. DISCUSSION
% ===================================================================
\section{Discussion}
\label{sec:discussion}

\subsection{What has been achieved}

\begin{enumerate}
\item \textbf{Reformulation.} $\PP$ vs $\NP$ is recast as an entropic No-Go theorem: polynomial-time algorithms cannot compress NP-witnesses below their Kolmogorov complexity.

\item \textbf{Barrier analysis.} The Kolmogorov-complexity approach satisfies all five barrier specifications: non-relativizing, non-natural, non-algebrizing, model-independent, and robust.

\item \textbf{Witness entropy gap.} There exist NP instances with witnesses of maximal conditional Kolmogorov complexity, while $\PP = \NP$ would force all witnesses to have logarithmic conditional complexity.

\item \textbf{Algorithmic positivity.} The concept identifies two core structural properties (polynomial-time witness extraction and entropic cheapness of witnesses) that $\PP$-decidable NP languages satisfy and NP-complete problems violate, supplemented by heuristic properties (compositional and monotone stability of witnesses) that provide additional separation evidence.
\end{enumerate}

\subsection{What remains open}

The central open problem -- the \textit{uniformity bridge} -- is the conversion of existential entropy bounds (``there exist hard instances'') to universal algorithm bounds (``no algorithm solves all instances''). This is the P-vs-NP-specific instance of the general pattern:

\begin{center}
\begin{tabular}{lll}
\hline
Problem & Positivity & Missing bridge \\
\hline
Hodge & HR-form $\geq 0$ & $\mathrm{AP} \Rightarrow$ algebraic \\
BSD & $\hat{h} \geq 0$ & Higher Gross-Zagier \\
P vs NP & $K^t(w|x)$ high & Uniformity bridge \\
\hline
\end{tabular}
\end{center}

\subsection{Limitations and intellectual honesty}

\textbf{This paper does not prove $\PP \neq \NP$.} We state this unambiguously. The contributions are: (1) a clean reformulation that identifies \textit{what} must be shown (witness incompressibility), (2) a barrier analysis arguing \textit{why} this approach is not blocked by known meta-theorems, and (3) concrete research directions (bridge strategies B1--B4) for closing the remaining gap.

\begin{enumerate}
\item \textbf{Non-computability of $K$.} Kolmogorov complexity is not computable, which means the entropy obstruction cannot be directly ``run'' as an algorithm. This is a feature (it bypasses natural proofs) but also a limitation (it makes the argument non-constructive).

\item \textbf{The unique-witness requirement.} The entropy gap argument is strongest for instances with unique witnesses. For instances with many witnesses, a polynomial-time algorithm might select a low-complexity witness even if high-complexity witnesses exist.

\item \textbf{Equivalence, not implication.} The Entropic Separation Conjecture is logically equivalent to $\PP \neq \NP$ (Remark~\ref{rem:esc-equivalence}). Therefore, this paper's No-Go theorem is a \textit{reformulation}, not an independent proof strategy. The value lies in identifying the information-theoretic structure of the problem.

\item \textbf{Barrier analysis is heuristic.} While we argue that $K$-based arguments are not naturally subject to the three barriers, the barrier theorems (relativization, natural proofs, algebrization) have precise technical definitions. A fully rigorous demonstration that the entropy approach formally escapes all three barriers would require showing that specific $K^t$-based statements cannot be relativized, which we leave as a research direction.
\end{enumerate}

\subsection{The meta-pattern}

Across all the problems in this research program, the structure is:

\begin{enumerate}
\item \textbf{Identify the positive form} (height, HR-form, capacity, entropy).
\item \textbf{Show one direction} (algebraic $\Rightarrow$ positive, $\PP \Rightarrow$ low entropy).
\item \textbf{Show the hard direction fails for obstructions} (non-algebraic violates positivity, NP-hard has high entropy).
\item \textbf{Close the loop with a composition/coupling theorem.}
\end{enumerate}

For P vs NP, steps 1--3 are established in this paper. Step 4 -- the uniformity bridge -- remains the central challenge.

\begin{remark}[Quantum extension]
\label{rem:quantum}
The entropy framework extends naturally to quantum computation. A quantum algorithm with $T$ gates on $n$ qubits can be described by $O(T \log T)$ classical bits (encoding the gate sequence), so $K(\mathrm{desc}(U)) \leq O(T \log T)$ for the classical description of the unitary $U$. If $\mathrm{BQP} = \mathrm{QMA}$, then the polynomial-time quantum algorithm would replace the QMA witness, and the analogous witness entropy gap argument would apply: the resource-bounded quantum Kolmogorov complexity of the witness would need to be simultaneously low (by the BQP algorithm) and high (for instances encoding random strings). This suggests $\mathrm{BQP} \neq \mathrm{QMA}$, though making this rigorous requires a quantum analogue of resource-bounded Kolmogorov complexity, which is an active area of research.
\end{remark}


% ===================================================================
% 12. CONCLUSION
% ===================================================================
\section{Conclusion}
\label{sec:conclusion}

The $\PP$ vs $\NP$ problem, like the Hodge Conjecture and the BSD conjecture, is fundamentally a statement about the \textit{impossibility of structural compression}: NP-witnesses cannot be algorithmically compressed to polynomial overhead, just as non-algebraic Hodge classes cannot be compressed to algebraic cycles, and analytic $L$-zeros cannot exist without geometric rational points.

The Kolmogorov complexity framework provides the natural positivity structure: it is non-negative, machine-invariant, non-computable, and captures the intrinsic information content of computational problems. The witness entropy gap -- the discrepancy between the high conditional complexity of witnesses and the low complexity that polynomial-time algorithms can produce -- is the fundamental obstruction.

The remaining challenge is the uniformity bridge: converting the existential entropy bound into a universal lower bound against all polynomial-time algorithms. This is the P-vs-NP-specific instance of the composition/coupling theorem that appears across all positivity-based approaches to fundamental mathematical problems.

\begin{quote}
\textit{``P vs NP does not ask whether hard problems exist. It asks whether the universe of efficient computation has room for problems whose solutions carry irreducible information -- and entropy theory provides the natural language to formulate this question with full precision.''}
\end{quote}


% ===================================================================
% ACKNOWLEDGMENTS
% ===================================================================
\begin{acknowledgments}
This work was conducted as independent research without institutional funding.

\textit{AI tools.} AI-based tools (Claude by Anthropic) were used for mathematical formalization and text generation. All content and responsibility for correctness lie solely with the author.

\textit{Funding information.} This research received no external funding.
\end{acknowledgments}


% ===================================================================
% BIBLIOGRAPHY
% ===================================================================
\begin{thebibliography}{30}

\bibitem{BakerGillSolovay1975}
T.~Baker, J.~Gill, \& R.~Solovay (1975).
Relativizations of the $\mathrm{P} =\mathrel{?} \mathrm{NP}$ question.
\textit{SIAM Journal on Computing}, 4(4), 431--442.
DOI: 10.1137/0204037.

\bibitem{RazborovRudich1997}
A.~Razborov \& S.~Rudich (1997).
Natural proofs.
\textit{Journal of Computer and System Sciences}, 55(1), 24--35.
DOI: 10.1006/jcss.1997.1494.

\bibitem{AaronsonWigderson2009}
S.~Aaronson \& A.~Wigderson (2009).
Algebrization: A new barrier in complexity theory.
\textit{ACM Transactions on Computation Theory}, 1(1), 2.
DOI: 10.1145/1490270.1490272.

\bibitem{LiVitanyi2008}
M.~Li \& P.~Vit\'anyi (2008).
\textit{An Introduction to Kolmogorov Complexity and Its Applications}, 3rd ed.
Springer.
DOI: 10.1007/978-0-387-49820-1.

\bibitem{Cook1975}
S.~A.~Cook (1975).
Feasibly constructive proofs and the propositional calculus.
\textit{Proceedings of the 7th ACM Symposium on Theory of Computing}, pp.\ 83--97.
DOI: 10.1145/800116.803756.

\bibitem{BenSasson2001}
E.~Ben-Sasson \& A.~Wigderson (2001).
Short proofs are narrow -- resolution made simple.
\textit{Journal of the ACM}, 48(2), 149--169.
DOI: 10.1145/375827.375835.

\bibitem{Razborov1985}
A.~Razborov (1985).
Lower bounds on the monotone complexity of some Boolean functions.
\textit{Doklady Akademii Nauk SSSR}, 281(4), 798--801.

\bibitem{Razborov1987}
A.~Razborov (1987).
Lower bounds on the size of bounded depth circuits over a complete basis with logical addition.
\textit{Mathematical Notes}, 41(4), 333--338.
DOI: 10.1007/BF01137685.

\bibitem{Grigoriev2001}
D.~Grigoriev (2001).
Linear lower bound on degrees of Positivstellensatz calculus proofs for the parity.
\textit{Theoretical Computer Science}, 259(1--2), 613--622.
DOI: 10.1016/S0304-3975(00)00403-7.

\bibitem{Mulmuley2011}
K.~D.~Mulmuley (2011).
On P vs.\ NP and geometric complexity theory.
\textit{Journal of the ACM}, 58(2), 5.
DOI: 10.1145/1667053.1667055.

\bibitem{Arora2009}
S.~Arora \& B.~Barak (2009).
\textit{Computational Complexity: A Modern Approach}.
Cambridge University Press.

\bibitem{Kolmogorov1965}
A.~N.~Kolmogorov (1965).
Three approaches to the quantitative definition of information.
\textit{Problems of Information Transmission}, 1(1), 1--7.

\bibitem{Chaitin1966}
G.~J.~Chaitin (1966).
On the length of programs for computing finite binary sequences.
\textit{Journal of the ACM}, 13(4), 547--569.
DOI: 10.1145/321356.321363.

\bibitem{Monasson1999}
R.~Monasson, R.~Zecchina, S.~Kirkpatrick, B.~Selman, \& L.~Troyansky (1999).
Determining computational complexity from characteristic `phase transitions'.
\textit{Nature}, 400, 133--137.
DOI: 10.1038/22055.

\bibitem{ImpagliazzoWigderson1997}
R.~Impagliazzo \& A.~Wigderson (1997).
P = BPP if E requires exponential circuits: Derandomizing the XOR Lemma.
\textit{Proceedings of the 29th ACM STOC}, pp.\ 220--229.
DOI: 10.1145/258533.258590.

\bibitem{GeigerFST2026}
L.~Geiger (2026).
The Renormalized Free Energy Framework.
Zenodo preprint.
\href{https://doi.org/10.5281/zenodo.19036191}{doi:10.5281/zenodo.19036191}.

\bibitem{GeigerRH2026}
L.~Geiger (2026).
The Riemann Hypothesis I--III.
Zenodo preprint.
\href{https://doi.org/10.5281/zenodo.19035845}{doi:10.5281/zenodo.19035845}.

\bibitem{Allender2006}
E.~Allender, H.~Buhrman, M.~Kouck\'y, D.~van~Melkebeek, \& D.~Ronneburger (2006).
Power from Random Strings.
\textit{SIAM Journal on Computing}, 35(6), 1467--1493.
DOI: 10.1137/050628994.

\bibitem{Hastad1987}
J.~H\r{a}stad (1987).
\textit{Computational Limitations of Small-Depth Circuits}.
MIT Press.

\bibitem{Cook1971}
S.~A.~Cook (1971).
The complexity of theorem-proving procedures.
\textit{Proceedings of the 3rd ACM Symposium on Theory of Computing}, pp.\ 151--158.
DOI: 10.1145/800157.805047.

\bibitem{ValiantVazirani1986}
L.~Valiant \& V.~Vazirani (1986).
NP is as easy as detecting unique solutions.
\textit{Theoretical Computer Science}, 47, 85--93.
DOI: 10.1016/0304-3975(86)90135-0.

\bibitem{Solomonoff1964}
R.~J.~Solomonoff (1964).
A formal theory of inductive inference, Parts I and II.
\textit{Information and Control}, 7(1), 1--22 and 7(2), 224--254.
DOI: 10.1016/S0019-9958(64)90131-7.

\bibitem{Allender2001}
E.~Allender (2001).
When worlds collide: Derandomization, lower bounds, and Kolmogorov complexity.
\textit{Proceedings of the 21st FSTTCS}, LNCS 2245, pp.\ 1--15.
DOI: 10.1007/3-540-45294-X\_1.

\bibitem{Hirahara2018}
S.~Hirahara (2018).
Non-black-box worst-case to average-case reductions within NP.
\textit{Proceedings of the 59th IEEE FOCS}, pp.\ 247--258.
DOI: 10.1109/FOCS.2018.00032.

\bibitem{Impagliazzo1995}
R.~Impagliazzo (1995).
A personal view of average-case complexity.
\textit{Proceedings of the 10th Structure in Complexity Theory Conference}, pp.\ 134--147.
DOI: 10.1109/SCT.1995.514853.

\end{thebibliography}

\end{document}
